\section{Resultados y Discusión}

\subsection{Presentación de Datos Experimentales}
% Presentar las mediciones cuantitativas en tablas y/o gráficas.
%
% Plantilla de tabla con booktabs:
% \begin{table}[htbp]
%     \centering
%     \caption{Mediciones experimentales.}
%     \label{tab:mediciones}
%     \begin{tabular}{cccc}
%         \toprule
%         \textbf{Ensayo} & \textbf{Variable $x$ [unidad]} & \textbf{Variable $y$ [unidad]} & \textbf{Incertidumbre} \\
%         \midrule
%         1 & 0.00 & 0.00 & $\pm 0.01$ \\
%         2 & 0.00 & 0.00 & $\pm 0.01$ \\
%         \bottomrule
%     \end{tabular}
% \end{table}

\subsection{Análisis Cuantitativo y Ajuste de Modelos}
% Describir regresiones, ajustes por mínimos cuadrados y estimación de parámetros con incertidumbres.

\subsection{Discusión de Resultados}
% Interpretar físicamente las tendencias observadas y evaluar concordancia con la teoría.